\section{Decision Analysis}
\label{sec:analysis}

R-MUR makes two decisions in sequence: which candidates enter the shortlist,
and which shortlisted candidate is selected. The analysis separates the
estimation error of the two stages, because the algorithm never evaluates the
response model outside the shortlist.

\subsection{Shortlist-aware regret}

Let $A$ be a fixed state, let
$j^\star=\argmax_{j\notin A}\marginal(j\mid A)$ be the best remaining
candidate, let
$j^\star_Q=\argmax_{j\in\mathcal Q_A}\marginal(j\mid A)$ be the best candidate
inside the shortlist, and let $\widehat{j}$ be the candidate the router selects
from the shortlist. The routing regret decomposes exactly into a screening term
and a reranking term,
\begin{equation}
 \underbrace{\marginal(j^\star\mid A)-\marginal(\widehat{j}\mid A)}_{\text{routing regret}}
 =
 \underbrace{\marginal(j^\star\mid A)-\marginal(j^\star_Q\mid A)}_{\text{screening regret}}
 +
 \underbrace{\marginal(j^\star_Q\mid A)-\marginal(\widehat{j}\mid A)}_{\text{reranking regret}} .
 \label{eq:decomposition}
\end{equation}
Each term is controlled by the error of the corresponding stage.

\begin{proposition}[Shortlist-aware one-step regret]
\label{prop:shortlist}
Suppose the cached score satisfies
$|\widehat m_C(j\mid A)-\marginal(j\mid A)|\le\epsilon_C$ for every remaining
candidate, and the response-aware score satisfies
$|\widehat m_R(j\mid A)-\marginal(j\mid A)|\le\epsilon_R$ for every
shortlisted candidate $j\in\mathcal Q_A$. Then the routing regret of
Eq.~\eqref{eq:decomposition} is at most
\begin{equation}
 \marginal(j^\star\mid A)-\marginal(\widehat{j}\mid A)\;\le\;2\epsilon_C+2\epsilon_R .
 \label{eq:shortlistbound}
\end{equation}
\end{proposition}

\begin{proof}
For the screening term, if $j^\star\in\mathcal Q_A$ the term is zero. Otherwise
every shortlisted candidate has a cached score at least that of $j^\star$,
because the shortlist consists of the $q$ largest cached scores. Hence
$\widehat m_C(j^\star_Q\mid A)\ge\widehat m_C(j^\star\mid A)$ and therefore
$\marginal(j^\star_Q\mid A)\ge\marginal(j^\star\mid A)-2\epsilon_C$. For the
reranking term, $j^\star_Q$ is feasible for the estimated argmax, so
$\widehat m_R(j^\star_Q\mid A)\le\widehat m_R(\widehat{j}\mid A)$, and two
applications of the response bound give
$\marginal(j^\star_Q\mid A)-\marginal(\widehat{j}\mid A)\le 2\epsilon_R$.
\end{proof}

The bound in Eq.~\eqref{eq:shortlistbound} is the statement that matches
Algorithm~\ref{alg:rmur}: the first term is the price of screening with cached
representations, the second is the price of estimating the response-aware score
on the retained set. The decomposition also exposes the cost--accuracy
trade-off of the shortlist: increasing $q$ reduces the empirical screening
regret and makes the best candidate more likely to reach the response-aware
stage, while requiring more predictor evaluations. Section~\ref{sec:experiments}
measures both sides of that trade-off.

\subsection{Stopping}

At deployment, R-MUR stops when no shortlisted candidate has a positive
estimated marginal utility. This stopping decision concerns the retained set:
candidates excluded by the cached screen are not evaluated by the response-aware
model. The shortlist-aware regret in Proposition~\ref{prop:shortlist}
quantifies the corresponding screening error, and the recall measurement in
Section~\ref{sec:experiments} reports how often the screen removes the best
candidate.

\subsection{Scope of the set-level statement}

The one-step results above require no assumption on the set function. A
guarantee for a sequence of $B$ selections would additionally require the
induced set utility to be monotone and submodular, so that approximate greedy
selection inherits the standard $(1-1/e)$ factor up to the accumulated
estimation error. Context utility is not monotone in general, because a context
can increase the loss, and we do not assume submodularity in the experiments.
We therefore state the sequential guarantee as an open condition rather than
claiming it for R-MUR.
